\section{Related Work}
\label{sec:related_work}

This review compares solar-forecasting studies by target, inputs, horizon,
and evaluation protocol. Section~\ref{sec:background} introduces the
underlying physical and architectural concepts.

\subsection{Forecasting scope and evaluation practice}

Solar-forecasting reviews organize methods by forecast horizon, input
source, target variable, and application
\cite{diagne2013,inman2013,antonanzas2016,voyant2017}. These distinctions
reflect different information sets and error structures: irradiance
nowcasting from sky images, hourly GHI forecasting from historical and
meteorological variables, day-ahead forecasting based on numerical weather
prediction, and monthly resource forecasting address different tasks.
Text-mining evidence also shows that the field has diversified in both
methods and applications \cite{yang2018}. Reviews focused on photovoltaic
power similarly find that reported performance depends on the horizon,
data, and optimization procedure \cite{ahmed2020}. This review therefore
treats GHI, solar irradiation, and photovoltaic power as distinct targets.

Evaluation practice is equally important. Deterministic solar forecasts
should be assessed against explicit baselines, with metrics selected for the
decision context and reported at the relevant horizons
\cite{yang2020verification}. General forecast-evaluation research also shows
that percentage measures and scale-dependent errors answer different
questions \cite{hyndman2006}. Voyant et al.\ \cite{voyant2022} demonstrate
that the appropriate naive baseline depends on temporal structure and
horizon. Their comparison of statistical baselines requiring no training
across multiple climates cautions against using a single persistence rule
as the sole benchmark in every setting.

\subsection{From intra-hour to week-ahead forecasting}

At very short horizons, cloud observations can provide information absent
from a univariate irradiance history. Chow et al.\ \cite{chow2011} derived
cloud maps and motion vectors from 30-s total-sky images at a distributed
test bed. They evaluated cloud conditions quantitatively and the derived GHI
nowcasts qualitatively. Reviews of intra-hour methods distinguish such
ground-imagery approaches from satellite, numerical weather prediction, and
time-series approaches \cite{chu2021,lin2023}. This literature also shows
why clear-sky or annual-average performance does not establish performance
during broken-cloud conditions.

For hourly time series, hybrid networks commonly combine representation
learning with recurrent dynamics. Huang et al.\ \cite{huang2021} used wavelet
packet decomposition, convolution, long short-term memory (LSTM), and
multilayer perceptron branches for 1-h-ahead prediction at three United
States sites. Their inputs included irradiance history, meteorological
variables, and next-hour forecasts of temperature, relative humidity, and
wind speed. Guariso et al.\ \cite{guariso2020} compared feed-forward and LSTM
networks using recursive, horizon-specific, and multi-output strategies for
horizons of up to 12~h. Their evaluation included persistence and clear-sky
baselines and an analysis of transfer across sites. These studies differ in
both model complexity and how they generate future steps.

Longer horizons can also retain hourly outputs. Cheng et al.\
\cite{cheng2021} combined one-dimensional convolution and LSTM components
for day-ahead to week-ahead daytime irradiance prediction, using recent
irradiance and a small set of weather variables. Zhao et al.\
\cite{zhao2024timesnet} coupled improved complete ensemble empirical mode
decomposition with adaptive noise (ICEEMDAN) to TimesNet, combining the
decomposed solar-radiation sequences with meteorological data for
week-ahead hourly forecasting. This application establishes the relevance
of TimesNet to solar forecasting, but its decomposition stage and
week-ahead protocol do not isolate the effect of the TimesNet backbone or
provide a matched comparison with DilatedRNN.

\subsection{Coarser resolution, multiple locations, and climate context}

Coarser targets lead to a different learning problem. Ozoegwu
\cite{ozoegwu2019} forecast monthly-mean daily global solar radiation with
nonlinear autoregressive, exogenous-input, and hybrid neural formulations.
Jung et al.\ \cite{jung2020} examined recurrent modeling for long-term
photovoltaic-power forecasting, whose target also depends on the conversion
system. Cabral et al.\ \cite{cabral2026} evaluated a GHI-only GPT-Neo
pipeline and contemporary competitors over 40 locations and a very-long-term
horizon. That work provides a precedent for evaluating long-horizon
forecasts across locations; the present study focuses on comparing two
time-series architectures with matched temporal resolution and forecast
design.

Cross-location learning can increase the effective sample size by sharing
parameters. It also changes the modeling setup from separate local models
to a shared forecasting model \cite{montero2021}. The broader recurrent
neural network (RNN) literature emphasizes the need to distinguish
architectural capacity from window construction, optimization, and
evaluation design \cite{hewamalage2021}. For seasonal and climate context,
Zhu et al.\ \cite{zhu2024} evaluated several deep time-series models across
horizons from minutes to 24~h and seasons, using different input combinations
and training-set lengths. Such designs
describe performance variation across conditions, but causal attribution
requires additional support from the study design.

The Supplementary Material summarizes these representative primary studies
in three groups organized by forecast horizon and geographic or climate
context. The tables compare dimensions relevant to this article's research
question, without ranking reported errors across studies with different
targets, sites, horizons, inputs, and test protocols.

\subsection{Research gap and positioning}

The reviewed studies provide evidence for specific architectures and
horizons, but do not establish how TimesNet and DilatedRNN compare under a
common protocol across temporal resolutions. Zhao et al.'s TimesNet
application predicts week-ahead hourly solar radiation after signal
decomposition \cite{zhao2024timesnet}; monthly neural studies use different
targets and architectures \cite{ozoegwu2019}; and the very-long-term
cross-location study centers on a language-model backbone
\cite{cabral2026}. Ranking their published errors would therefore confound
model choice with temporal aggregation, information availability, location
set, and evaluation protocol.

The present study addresses this comparison through a common experimental
design. At each resolution, TimesNet and DilatedRNN share the same samples,
chronological forecast origins, admissible inputs, post-processing rules,
baselines, and aggregation across locations. Evaluation distinguishes hourly
multi-day forecasts from daily aggregates and treats the primary one-step
monthly task separately from the exploratory six-month task. These choices
address the comparability limitations identified in solar-forecast
verification and benchmarking research
\cite{yang2020verification,voyant2022}.
